\chapter{Summary and Conclusions}\label{ch:7}

\section{Conclusions}\label{sec:7.1}

All results below come from models trained until their training accuracy reached (nearly) 100 \%.

\begin{enumerate}
  \item \textbf{Best on all patients: 0.558} [0.530, 0.587] from image-derived inputs only, by a ten-member ensemble with an ordinal meta-learner that includes EchoNext-trained members. The best single model (E2 LightGBM, 0.557) and the best image-only ensemble (ENH-vote, 0.549, no printed measurements) are close behind.
  \item \textbf{Best in the 60s: 0.696} [0.630, 0.762] on the 48 \% of patients the image-only triage model answers; 89 \% of its confident ``Normal'' answers have an echo EF of at least 41 \%.
  \item \textbf{Clean labels are graded well:} 0.67--0.69 on clear-cut and noise-free subsets (0.694 on 264 clean ECGs).
  \item \textbf{Transfer works across populations:} models trained only on US public data reach up to 0.539 four-class accuracy on our printouts with only their thresholds tuned, and add about two points to the best ensemble.
  \item \textbf{Fitting the training data perfectly does not carry over to new patients:} every flexible model reached 1.000 training accuracy, while held-out accuracy stayed at 0.35--0.56 on all patients.
  \item \textbf{The ceiling is set by the labels and the information in the ECG:} 70 \% on all patients would need EF estimates correlating at $r\approx0.90$ with echo, against 0.62 for our best EF regressor.
\end{enumerate}
\section{Scope for future work (Phase II)}\label{sec:7.2}

\begin{itemize}
  \item Larger, multi-year data with biplane-Simpson EF linked by patient identifier within days of the ECG.
  \item Digital 10-s simultaneous ECGs where available, and enrichment for Mild and Moderate patients.
  \item Prospective evaluation of the triage model as an echocardiography-referral aid, and deployment as a lightweight application.
\end{itemize}
